\section{Background}
\label{sec:background}

\subsection{FragPicker's analysis phase}
FragPicker~\cite{park2021fragpicker} observes that on flash and Optane devices the
harm done by fragmentation comes mostly from request splitting, so it moves only
the pieces of a file that the application's own requests straddle. Its public
artifact~\cite{fragpickerartifact} implements this in two phases. The
\emph{analysis phase} is a chain of scripts in \texttt{src/analysis}:
\begin{enumerate}
\item \texttt{trace.sh} starts two BCC~\cite{bcc} \texttt{trace} processes, one
  on \texttt{vfs\_read} and one on \texttt{vfs\_write}, filtered by process name.
  Each event prints the inode number, request size, file position, whether the file
  was opened with \texttt{O\_DIRECT}, and the request type. Both processes append to
  \texttt{./trace.result}.
\item \texttt{parse.sh} keeps those five fields and sorts the lines by inode
  number (stable).
\item \texttt{processing.py} has two loops. Phase~1 aligns every request to 4\,KiB
  blocks, models the kernel's 128\,KiB readahead window for sequential buffered
  reads, and appends one line per request to \texttt{<inode>.txt}; it opens and
  closes that file once per trace line. Phase~2 visits every inode seen in the
  trace, runs \texttt{find /mnt -inum <inode>} to turn it into a path, drops
  directories and inodes that \texttt{find} cannot see, sorts the file with
  \texttt{sort -g} and \texttt{mv}, and appends ``\texttt{<inode> <count>}'' to
  \texttt{filelist.txt}.
\item \texttt{merge.py} merges overlapping requests of each file into windows with
  a hit count (one \texttt{mv} subprocess per file).
\item \texttt{hotness.sh} keeps the hottest $p\%$ of the windows of each file with a
  shell loop of \texttt{wc}, \texttt{awk}, \texttt{sort} and \texttt{head}.
\end{enumerate}
The \emph{migration phase} (\texttt{src/migration}) then resolves each listed inode
with \texttt{find} again, reads the file's extents with \texttt{filefrag -v}, and
rewrites every hot window that crosses an extent boundary, after punching and
re-allocating its blocks on in-place file systems such as ext4.

\begin{table}[t]
\caption{Work per traced file in each analysis step, counted with \texttt{strace -f}
on 50 traced files (E1c). Upstream's per-file \texttt{find} also reads every
directory entry under \texttt{/mnt}.}
\label{tab:anatomy}
\small
\begin{tabular}{@{}lrrr@{}}
\toprule
Step & processes & \texttt{openat} & per tree entry \\
\midrule
\texttt{processing.py} (upstream) & 3 & 117 & readdir, $\times F$ \\
\texttt{merge.py} (upstream) & 1 & 39 & -- \\
\texttt{hotness.sh} (upstream) & 10 & 174 & -- \\
rewrite, all three steps & 0 & 4.7 & readdir, $\times 1$ \\
\bottomrule
\end{tabular}
\end{table}

Table~\ref{tab:anatomy} counts what these steps cost per traced file (totals divided
by 50, so interpreter start-up is included): fourteen new processes, one of which
walks the whole tree, and about 330 \texttt{openat} calls, most of them made by the
new processes while they load their libraries.

Two properties of this design matter below. First, the analysis phase never looks
at file layout: extents are read only during migration. Second, the artifact's own
evaluation traces one file per experiment (\texttt{evaluation/read\_benchmark}), so
the per-file loop of Phase~2 runs a handful of times; how it scales with the number
of traced files was never part of the evaluation.

\subsection{The parallel analysis fork}
A later fork~\cite{forkbranch} changes Phase~2 behind a \texttt{--parallel} flag. It
builds one inode-to-path dictionary with a single
\texttt{find -type f -printf '\%i \%p\textbackslash n'}, then gives each traced
inode to a pool of Python threads. Each worker looks the path up in the dictionary,
calls \texttt{stat}, and reads the file's extents with the FIEMAP
ioctl~\cite{kernelfiemap} (\texttt{stat}, \texttt{open}, one or more
\texttt{ioctl} calls, \texttt{close}). The extents of every file are then written to a
temporary directory and sorted there with \texttt{sort} subprocesses; a third step
sorts every \texttt{*.txt} file in the working directory with more \texttt{sort}
subprocesses; and \texttt{filelist.txt} lists the inodes whose FIEMAP call succeeded.
Phase~1 is unchanged. The fork reports a 12.74$\times$ speedup with 16 threads for
1{,}000 files of 100\,KB when a 5\,ms \texttt{time.sleep} is added before every
extent lookup, and a slowdown (0.10--0.26$\times$) without the delay on an NVMe SSD.

\subsection{What a per-file metadata lookup costs}
On ext4 the root of a file's extent tree lives in the inode itself and holds up to
four extents; larger trees add index and leaf blocks~\cite{kernelext4}. An inode
table block holds several inodes (sixteen with the default 256-byte inode), and when
ext4 has to read an inode table block it reads ahead up to 32 neighbouring blocks by
default (\texttt{inode\_readahead\_blks}). Files created together therefore tend to
share the device reads that bring their inodes into memory. A lookup of a cached
inode with an inline extent tree needs no device I/O at all. A fixed per-file delay
models none of this: it charges the same latency to every lookup, cached or not,
and lets any number of lookups wait at the same time.
